\subsubsection{Energia cinética}
\label{sec:dinAcoplada_energia_cinetica}

A energia cinética total é escrita como a soma das contribuições da base e dos $n_\ell$ elos do manipulador robótico

\begin{equation}
T = T_{B} + \sum_{i=1}^{n_\ell} T_{i},
\label{eq:T_total}
\end{equation}

onde $T_B$ é a energia cinética translacional e rotacional da base e $T_i$ é a energia cinética do $i$-ésimo elo, expressas no referencial da base $\{B\}$, com $n_\ell = 3$ para o \gls{mr} deste trabalho.
A energia cinética da base é escrita, no caso em que apenas a atitude é considerada como grau de liberdade dinâmico, como

\begin{equation}
T_{B} =
\frac{1}{2}\, {}^{B}\!\vec\omega_{B}^{\,T}\, I_B\, {}^{B}\!\vec\omega_{B},
\label{eq:TB}
\end{equation}

em que $\mathbf{I}_B$ é o tensor de inércia da base em torno do seu centro de massa, expresso no referencial $\{B\}$.
No modelo adotado, a base física não possui graus de liberdade translacionais na formulação acoplada, uma vez que a translação é descrita separadamente pelo movimento orbital relativo em \gls{lvlh}.
Então, apenas o termo rotacional associado à velocidade angular ${}^{B}\!\vec\omega_B$ contribui para a energia cinética e para a matriz de inércia generalizada $M(\vec q)$.

\begin{equation}
T_i =
\frac{1}{2} m_i\, {}^{B}\!\vec v_{G_i}^{\,T} {}^{B}\!\vec v_{G_i}
+
\frac{1}{2} {}^{B}\!\vec\omega_i^{\,T} I_i {}^{B}\!\vec\omega_i,
\label{eq:Ti}
\end{equation}

onde ${}^{B}\!\vec v_{G_i}$ e ${}^{B}\!\vec\omega_i$ são a velocidade linear do centro de massa e a velocidade angular do elo $i$, obtidas a partir do Jacobiano acoplado da Seção~\ref{sec:cinematica_diferencial_jacobiano}.

Agrupando-se as velocidades generalizadas no vetor

\begin{equation}
\dot{\vec y}
=
\begin{bmatrix}
{}^{B}\!\vec\omega_{B} \\
\dot{\vec q}
\end{bmatrix}
\in \mathbb{R}^{6}, 
\end{equation}

a energia cinética total pode ser escrita de forma compacta como

\begin{equation}
T =
\frac{1}{2}\,\dot{\vec y}^{\,T}\,
M(\vec q)\,
\dot{\vec y},
\label{eq:T_quadratica_Mq}
\end{equation}

em que $M(\vec q)\in\mathbb{R}^{6\times 6}$ é a matriz de inércia generalizada do conjunto base--manipulador, construída a partir das contribuições da base e dos elos segundo as expressões \eqref{eq:TB} e \eqref{eq:Ti}.
